\section{Related Work}\label{sec:related}

\subsection{Selfish routing and the price of anarchy}
The study of routing as a game predates learning-based methods. Wardrop's principles
distinguish the user equilibrium, where every used route between an origin and
destination has equal and minimal cost to the individual, from the system optimum that
minimizes total cost~\citep{wardrop1952}. \citet{roughgarden2002selfish} quantified how far apart these can be, proving
bounds on the price of anarchy for selfish routing, and the Braess paradox shows the equilibrium can
be counter-intuitively bad~\citep{braess2005paradox}. The gap is not merely theoretical:
\citet{zhang2018poa} estimated the price of anarchy from real traffic data on
the Eastern-Massachusetts road network and proposed data-driven strategies for reducing
it, confirming that deployed networks operate measurably above the system optimum. These
results motivate selfless routing directly: the potential benefit of asking vehicles to sacrifice is exactly the
price-of-anarchy gap, which is large only when the network forces contention onto scarce
links. Our experiments make this dependence concrete: selfless detours help precisely
in the saturated scenarios where the gap is wide, and not otherwise.

\subsection{The selfless traffic routing model}
\citet{dai2019str} introduced the selfless driving model and argued that
vehicular routing should be evaluated by network-level average travel time rather than
individual shortest paths. Their formulation permits a controlled vehicle to take a locally
suboptimal route when doing so benefits the group. We adopt this objective but replace
their heuristic controller with a learned cooperative policy, and we ground the trade-off
in a common travel-time unit so the sacrifice decision is made on principled terms.

\subsection{Reinforcement learning and MAPPO}
Reinforcement learning is a natural fit for sequential routing, where actions change the
state that later actions face~\citep{mnih2015human}. Proximal Policy Optimization (PPO)
provides a stable on-policy actor--critic update via a clipped surrogate
objective~\citep{schulman2017ppo}, and Generalized Advantage Estimation (GAE) reduces the
variance of the advantage signal it optimizes~\citep{schulman2016gae}. In the multi-agent
setting, the environment becomes non-stationary from any single agent's view. The
centralized-training / decentralized-execution paradigm addresses this by giving the
critic access to global information during training while the actor remains
local~\citep{lowe2017maddpg}. MAPPO applies PPO under CTDE with a centralized value
function and has proven surprisingly strong on cooperative
benchmarks~\citep{yu2022mappo}; we use it here because our reward is cooperative and our
deployment must be decentralized (each vehicle routes from its own view).

\subsection{Credit assignment in cooperative MARL}
A cooperative reward shared identically by all agents makes it hard to tell which agent's
action helped, because each agent's contribution is buried in a large common signal.
Difference rewards address this by crediting each agent with its marginal effect on the
global objective~\citep{wolpert2001collectives}, and counterfactual multi-agent policy
gradients (COMA) estimate that marginal effect with a centralized
critic~\citep{foerster2018coma}. Our team-reward term is a leave-one-out difference reward
in travel-time units: each agent pays for its own slowness relative to the live fleet
average, so the external, demand-driven congestion level, which no single vehicle
controls, cancels out and the learnable signal is sharpened.

\subsection{Simulation environment}
We use Eclipse SUMO as the microscopic (vehicle-level) traffic simulator~\citep{lopez2018sumo,krajzewicz2012recent}.
SUMO models car-following, lane changing, junctions, and traffic lights, which makes the
evaluation considerably more realistic than routing on a static weighted graph, but also
makes control harder: a route that looks good at the graph level is only useful if the
vehicle can physically execute it under lane-feasibility, route-continuity, and
commit-window constraints. Dijkstra's algorithm~\citep{dijkstra1959note} remains the
appropriate baseline, not as a weak straw man but as a strong, deterministic controller
that respects the road graph and never over-commits, which is why it is competitive in
exactly the low-congestion regime where selfless routing has nothing to gain.
